\section{おわりに}

本実験では，昼食後のデスクワークにおける居眠りの防止を目的として，Arduino と光センサを用いた居眠り検知・防止装置を開発した．脈拍センサで取得した値を用いてLPとF-stressを表示し，被験者のストレス状態を幾何学的かつ数値的に扱った．それらを用いて本装置の居眠り検知・防止を検証したところ，本装置で居眠り検知・防止が可能であることが明らかになった．また，誤使用時や暗所時におけるシステムの動作の確認を行い，正しく実装できていることがわかった．システム側から誤使用を検知でき，既存の評価指標ではできなかったうとうとしている様子をグラフ化できる点において，本システムは有効性がある．しかし，うとうとしている時の頭の動きをグラフ上に表すことが出来たが，プログラム側でうとうと状態を判定するシステムは未実装である．うとうとしている状態をシステム側が検知できれば，深い居眠りに到達する前に被験者を起床させることができる可能性がある．その実装のためには，うとうとしている時の光センサの値とその時のLPのデータをより多く収集し，波形を解析する必要がある．